\section{Discussion}
\label{sec:final}

We investigated a fragment of the query language XPath from a Description Logic (DL) perspective. 
The fragment under consideration (called $\xpath$) features data comparison and is interpreted over data graphs where accessibility relations are functional. We showed that even with this restriction we still model interesting scenarios. 

The navigational fragment of XPath has been studied in a DL setting before~\cite{KostylevRV14,Kost:xpat15}, but to our knowledge this has not been done for fragments featauring data comparisons. We start by presenting a novel connection between $\xpath$ and the DL with concrete domains $\alcd$. This enables us not only to see $\xpath$ as a DL with a very simple concrete domain, but also to obtain complexity upper bounds for reasoning taks. Moreover, we showed that this bounds can be refined and thus $\xpath$ has good computational properties. Concretely, we proved that concept satisfiability, ABox consistency and concept satisfiability w.r.t. a TBox are all \np-complete problems. This is is contrast with DL with concrete domains in general, where the complexity is usually quite high. 

There are still several lines of work to be explored. Firstly, it would be interesting to further explore the relational version of XPath (called here $\xpathr$) in a DL setting, and obtain tight complexity bounds. Secondly, we would like to extend our methods and results to related logics, as e.g. fragments of the SHACL language~\cite{Ortiz23,Ahmetaj0S23}. Finally, we would like to implement XPath reasoners for the tasks we investigated.